% Auto-generated by experiments/19_report_visualization_enrichment.py
\section*{Visual Evidence Gallery}
\begin{figure}[!htbp]
\centering
\begin{subfigure}[t]{0.48\linewidth}
\centering
\includegraphics[width=\linewidth]{results/figures/report_enrichment/stage01_math_notes_payoff_curves.pdf}
\caption{Put and call payoff curves establish the obstacle that all American option values must respect.}
\end{subfigure}
\hfill
\begin{subfigure}[t]{0.48\linewidth}
\centering
\includegraphics[width=\linewidth]{results/figures/report_enrichment/stage01_math_notes_exercise_continuation_regions.pdf}
\caption{The exercise and continuation regions show why the early-exercise boundary is a first-class diagnostic rather than a decorative price threshold.}
\end{subfigure}
\end{figure}

\begin{figure}[!htbp]
\centering
\begin{subfigure}[t]{0.48\linewidth}
\centering
\includegraphics[width=\linewidth]{results/figures/report_enrichment/stage01_math_notes_complementarity_geometry.pdf}
\caption{The variational inequality couples obstacle contact with zero PDE residual in continuation.}
\end{subfigure}
\hfill
\begin{subfigure}[t]{0.48\linewidth}
\centering
\includegraphics[width=\linewidth]{results/figures/report_enrichment/stage01_math_notes_continuation_premium.pdf}
\caption{The continuation premium C equals U minus Phi and separates immediate-exercise payoff from the value of waiting; this motivates the later price-only surrogate.}
\end{subfigure}
\end{figure}

\begin{figure}[!htbp]
\centering
\begin{subfigure}[t]{0.48\linewidth}
\centering
\includegraphics[width=\linewidth]{results/figures/report_enrichment/stage01_math_notes_spot_grid_stencil.pdf}
\caption{Three neighboring spot nodes show how first and second derivatives are approximated on the grid.}
\end{subfigure}
\hfill
\begin{subfigure}[t]{0.48\linewidth}
\centering
\includegraphics[width=\linewidth]{results/figures/report_enrichment/stage01_math_notes_coefficient_decomposition.pdf}
\caption{The lower, diagonal, and upper coefficients show how diffusion, drift, and discounting enter one row of the finite-difference operator.}
\end{subfigure}
\end{figure}

\begin{figure}[!htbp]
\centering
\begin{subfigure}[t]{0.48\linewidth}
\centering
\includegraphics[width=\linewidth]{results/figures/report_enrichment/stage01_math_notes_interior_vector_layout.pdf}
\caption{The solver matrix acts on interior unknowns while boundary values enter through the right hand side.}
\end{subfigure}
\hfill
\begin{subfigure}[t]{0.48\linewidth}
\centering
\includegraphics[width=\linewidth]{results/figures/report_enrichment/stage01_math_notes_annotated_operator_matrix.pdf}
\caption{The tridiagonal matrix layout maps the stencil coefficients into rows of the discrete operator.}
\end{subfigure}
\end{figure}

\begin{figure}[!htbp]
\centering
\begin{subfigure}[t]{0.48\linewidth}
\centering
\includegraphics[width=\linewidth]{results/figures/report_enrichment/stage01_math_notes_cn_matrix_pair.pdf}
\caption{The Crank--Nicolson matrices A and B split the same spatial operator across the new and old time levels; the diagram explains construction rather than adding accuracy evidence.}
\end{subfigure}
\hfill
\begin{subfigure}[t]{0.48\linewidth}
\centering
\includegraphics[width=\linewidth]{results/figures/report_enrichment/stage01_math_notes_lcp_vector_form.pdf}
\caption{The vector inequalities connect the payoff obstacle with the matrix residual at every interior node.}
\end{subfigure}
\end{figure}

\begin{figure}[!htbp]
\centering
\begin{subfigure}[t]{0.48\linewidth}
\centering
\includegraphics[width=\linewidth]{results/figures/report_enrichment/stage01_math_notes_time_marching_grid.pdf}
\caption{The finite difference method starts from the payoff and advances through time to maturity.}
\end{subfigure}
\hfill
\begin{subfigure}[t]{0.48\linewidth}
\centering
\includegraphics[width=\linewidth]{results/figures/report_enrichment/stage01_math_notes_rhs_boundary_adjustment.pdf}
\caption{Boundary values affect only the edge equations of the interior linear system.}
\end{subfigure}
\end{figure}

\FloatBarrier
